\section{可直接执行的 Codex 任务书}
完整可复制文本随交付文件 \texttt{Codex\_R7\_Minimal\_Experiment\_Prompt.txt} 提供。默认 MODE=A，执行 E1+E2，总上限 50,000 complex RHS；MODE=B 为独立后续任务，不自动启动。两种模式的目的、对象、四臂、成本、reference、统计和失败分类均已明确，避免把理论研究变成所有方向一起运行的 campaign。

\subsection{默认任务的操作摘要}
\begin{enumerate}
\item 固定科学提交，在独立目录恢复真实 P1 资产。不能用简化 forward 冒充 Full Maxwell，不能覆盖历史账本。
\item 做 256 新 decoder 点的零物理 gauge 审计；generic 满秩时删除普遍降维主张，但继续 E2。
\item 固定 4 开发和 8 locked 新中心、两个采集、两个方向、四距离；按冻结 prior 种子顺序抽取，不因秩或效果重抽。局部几何不要求 inverse MAP 驻点。
\item 实现 material-regular current jets、明确强补偿类、误差及有限有效域。满秩点 exact $p=1$，不得人为抬阶。
\item 对四臂作同材料距离、同总费用比较；保存完整响应/current/系数向量，不只存 norm。
\item 用共同低阶材料系数计算采集补偿冲突，和 joint Fisher、secant、真实 nonlinear profile 比较。
\item locked 预测与费用先封存，再评分。按 parent 统计，不把源、方向、seed 当独立对象。
\item 完成之后逐 claim 判 T/I/B/P/C；保留正面材料探索事实与具体失败，不自动追加大型网络。
\end{enumerate}

\subsection{必须保留的执行约束}
用户给定原 prior、noise 和 solver 不变；不得为了出现多解放宽 noise-ball 或偷换材料度量。相容集合不是 posterior 硬支持。已知合法几何中心不冒称未知重建。若资产不足，完成可做的推导/实现并说明缺口；如果只是代码或环境错误，在授权预算内修复继续。物理账本中央唯一；所有失败与重试计费。任何缺乏严格误差常数的结果用“经验收敛”或“有限尺度支持”表述。

\section{十五个核心问题的明确回答}
\label{sec:answers}
\begin{enumerate}[label=\textbf{\arabic*.},leftmargin=2em]
\item \textbf{从多解探索到任务计算，是深化还是漂移？}\\
材料商与概率体积是深化；任务统计是有独立风险证据支持的自然分叉；缩短 Native 当前主要是工程优化，未被前两者推出。不能把整个演化全部合理化，也不能全部否定。

\item \textbf{Maxwell 材料多解的正确数学对象是什么？}\\
带实际材料 prior 的可实现材料集合及 $F$ 的 exact fibers/finite-noise likelihood tubes；概率对象为 $L_y\,d\mu_0/Z_y$，可含原子、不同维 strata 和重叠原像。单张光滑流形只在适当常秩区域成立。

\item \textbf{Dynamic weak directions 的原创价值还剩多少？}\\
作为泛化算法原则原创性有限；作为材料正则的有限幅度候选生成器，以及识别补偿、主导区间和采集冲突的入口，仍有实质价值。非对合现象不能单独成为论文中心。

\item \textbf{如何区分 gauge 与真实 Maxwell 不可辨识？}\\
先验证材料差是否非零，再用材料度量分析 $DG$ 与 $DF$；有限等材料见证定义 gauge，单点 $DGv=0$ 不够。必须记录材料变化、响应变化、可行 prior 与数值误差四者。

\item \textbf{应否直接在材料商空间 Bayesian inference？}\\
统计目标应以材料测度定义。计算上，在 exact gauge 可积分且商核便宜的区域可直接做；generic 满秩、chart 复杂或核难算时保留 latent 并维护正确联合条件分布更好。不能在混合维商上强塞一个全维 Lebesgue Flow。

\item \textbf{Full66 posterior 为何未合格？}\\
已观察到高 SNR 下 importance overlap/ESS、链混合、强弱条件近似、绝对 response calibration 与预算问题。R6 又揭示条件依赖和材料奇性，但尚未证明哪一项是全部失败的唯一根因。有限 Gaussian likelihood 下 latent posterior 仍对原 Gaussianized prior 绝对连续；Flow 不是因材料原子而数学上不可能。方法仍需更好的条件结构、覆盖与成本设计。

\item \textbf{什么概率体积真正必要？}\\
正确材料 prior 权重、影响指定任务的分支质量，以及遗漏区域的任务贡献界。只求材料任务可积分掉 decoder fiber；要模式概率、事件、主动测量或 latent 事件时不能删除相应核与质量。不要恢复所有 latent 密度，也不要任意删体积。

\item \textbf{rank、effective rank、conditioning 与任务信息是什么关系？}\\
rank 是局部代数维数；effective rank 相对 prior、noise、可行幅度；conditioning 量化局部反演稳定性；任务信息还取决于任务函数与全局分支概率。正则单峰线性近似时可相互联系，临界、对称、多模态及分层时差别很大。

\item \textbf{是否有超越局部 SVD 的 task-conditional observability？}\\
有：完整条件期望信息算子 $\mathcal O_{a,y}$、有限任务风险下降、材料 secant/补偿结构。它们不是全新的一般 Bayesian 原理。$B_\varphi$ 仅给一阶响应，有明确高阶盲点；Maxwell 原创性需体现在其可计算物理结构。

\item \textbf{能否由任务决定 Jacobian/Maxwell 精度？}\\
在局部模型或有可靠状态、积分、尾部界时可以。式 \eqref{eq:task-cov-error} 无需稳定每个弱特征向量；式 \eqref{eq:centered-certificate} 可给任务停止。当前 Full66 全域廉价证书没有建立，不能把学习 risk predictor 当作它。

\item \textbf{R6 加速是否只是修复原流程？}\\
现有证据更支持这种解释：两阶段坐标 MAP 不同，Direct Probit 强，Full-q 控制缩小优势。它仍是有价值的实际计算改进；没有独立机制之前，不以新 Maxwell 定理包装。

\item \textbf{哪条线最可能产生新数学定理？}\\
真实材料补偿的有限有效域、共享采集补偿冲突及其对任务可区分性的界。高阶递推和最小二乘恒等式自身是先例；需要把可验证常数、物理约束和 finite-scale 用途连起来。

\item \textbf{哪条线较快完成有说服力算法论文？}\\
材料任务成像 B 已有最成熟独立风险证据；补上物理特征的公平增量与任务资格较快。若 E2 证明冲突指标能替代昂贵 nonlinear profile，则 A 的聚焦采集/探索算法也可能更短。现在都不能承诺强 TAP 已具备。

\item \textbf{哪些永久放弃，哪些暂停？}\\
永久放弃错误推论，见原创性节；暂停具体无增量的大训练、无权重轨迹平均与过度阶段优化。真实多解、正确概率输运、材料函数量与动态弱几何仍值得研究。

\item \textbf{写一篇、两篇还是先证明核心？}\\
先证明并检验一个核心：材料补偿为何有效、何时失效、怎样被新采集破坏。随后 A/B 分成两篇或以自然的共同算法合并；C 暂为模块。当前不写覆盖所有 R2--R6 的统一大论文。
\end{enumerate}

\section{最终科学陈述}
\begin{tcolorbox}
有限观测约束的首先是一类材料的电磁响应，而不是一幅唯一材料图像。同一种材料可以有多种参数表示；不同材料又可以通过共同的物理补偿产生几乎相同的观测。研究的关键是，在正确材料测度上分清这两种不确定性，确定真实材料差异在什么尺度、哪种采集下才开始可见，并理解这些差异如何改变指定材料判断。由此产生的新科学问题是：\textbf{材料间可补偿的响应关系，如何决定观测能够消除多少任务相关不确定性，以及为了达到该任务精度究竟需要计算哪些物理信息。} 这条关系有明确的 Maxwell 机制和可证伪的数值入口；其完整理论与算法优势仍需下一轮集中验证。
\end{tcolorbox}

\subsection{交付与可复现说明}
本报告不重写 R2--R6 的历史资格；所有新数值来自既有封存数组复算或明确标注的解析示意，不包含新增 Maxwell 实验。完整包提供 PDF、可编辑 LaTeX 源、可复制 Codex 任务书、来源索引和审查记录。正文中的命题区分了成熟数学应用、Maxwell 特化和待证的论文贡献；未将未来实验、推测的原创性或工程节时写成既成理论结果。

